\documentclass[11pt]{article}
\usepackage[margin=1in]{geometry}
\usepackage{booktabs,graphicx,hyperref,xcolor,amsmath}
\usepackage[numbers]{natbib}
\newcommand{\todo}[1]{\textcolor{red}{[TODO: #1]}}
\newcommand{\tbd}{\textcolor{blue}{[TBD]}}

\title{Does Multi-Agent Decomposition Increase or Reduce Prompt-Injection Risk?\\
\large The Role of Handoff Format and User-Intent Preservation}
\author{\todo{Member 1} \and \todo{Member 2} \and \todo{Member 3} \and \todo{Member 4}}
\date{November 2026}

\begin{document}
\maketitle

\begin{abstract}
\todo{~150 words. (1) Context: multi-agent workflows split reading, planning and execution.
(2) Gap: which handoff design choices govern injection propagation. (3) Method: single agent
vs.\ four reader--planner--executor pipelines in a $2\times2$ design (free-text vs.\ JSON handoff
$\times$ executor sees the user's request), two attack goals, \tbd{} episodes.
(4) Results: \tbd. (5) Takeaway: \tbd.}
\end{abstract}

% ---------------------------------------------------------------- 0.75 p
\section{Introduction}
\todo{P1 motivation: agents act on untrusted content; indirect prompt injection~\cite{greshake2023}.}
\todo{P2 two competing views: isolation vs.\ laundering. Point to Fig.~\ref{fig:pipeline}.}
\todo{P3 our question: which handoff design choices change risk, not ``more vs.\ fewer agents''.}
Our contributions are:
\begin{itemize}
  \item a controlled $2\times2$ comparison (handoff format $\times$ user-intent preservation) plus a single-agent baseline;
  \item joint measurement of redirection, content manipulation, task utility and cost;
  \item trace-level and replay analysis locating where injections are dropped, quoted or promoted;
  \item released code, scenarios and traces.
\end{itemize}

\begin{figure}[t]
  \centering
  \fbox{\parbox{0.9\linewidth}{\centering\vspace{3em}\todo{Pipeline diagram with a laundered instruction crossing a handoff}\vspace{3em}}}
  \caption{Reader $\rightarrow$ planner $\rightarrow$ executor workflow and the two manipulated design variables.}
  \label{fig:pipeline}
\end{figure}

% ---------------------------------------------------------------- 1 p
\section{Background and Related Work}
\paragraph{Prompt injection in tool-using agents.} \todo{\cite{greshake2023,injecagent2024,agentdojo2024}}
\paragraph{Injection in multi-agent systems.} \todo{\cite{promptinfection2024,triedman2025,pear2026,planflip2026}}
\paragraph{Architectural defenses.} \todo{dual LLM~\cite{willison2023dual}, design patterns~\cite{beurerkellner2025patterns}, CaMeL~\cite{camel2025}, multi-CaMeL~\cite{camelstalk2026}: state exactly how we differ.}
\paragraph{Prompt-level defenses.} \todo{spotlighting~\cite{spotlighting2024}, LLM tagging~\cite{promptinfection2024}; relation to our replay study.}
\paragraph{Positioning.} We do not claim to discover cross-agent injection or laundering; we isolate two design variables within one fixed workflow.

% ---------------------------------------------------------------- 0.75 p
\section{Threat Model and Task}
\todo{Task; attacker capability (one document); G1 redirect; G2 manipulation and its boundary (quoting inside a warning is not success); out of scope.}
\begin{table}[h]\centering\small
\begin{tabular}{@{}ll@{}}\toprule
Element & Definition\\\midrule
Victim task & Summarize 3--4 meeting documents, email the user-named recipient(s)\\
Attacker & Edits one document; cannot touch prompts, tools or inter-agent messages\\
G1 & Email sent to the attacker's address (to/cc/bcc); other unauthorized recipients are scored as misdirection\\
G2 & Attacker URL presented as normal content to the authorized recipient\\
Knowledge & Generic templates vs.\ architecture-aware (\texttt{role}) template\\\bottomrule
\end{tabular}
\caption{Threat model.}\label{tab:threat}
\end{table}

% ---------------------------------------------------------------- 1.25 p
\section{System Design}
\todo{Tools (mocked); roles and permissions; handoff formats and schemas (App.~\ref{app:schemas}); repair rule; executor intent view; controls; optional gateway.}
\begin{table}[h]\centering\small
\begin{tabular}{@{}lllc@{}}\toprule
ID & Architecture & Handoff & Executor sees request\\\midrule
A  & single agent & --- & yes\\
M1 & reader$\to$planner$\to$executor & free text & no\\
M2 & reader$\to$planner$\to$executor & JSON & no\\
M3 & reader$\to$planner$\to$executor & free text & yes\\
M4 & reader$\to$planner$\to$executor & JSON & yes\\\bottomrule
\end{tabular}
\caption{Configurations.}\label{tab:configs}
\end{table}

% ---------------------------------------------------------------- 1 p
\section{Experimental Setup}
\todo{Scenarios (\tbd{} test, 5 dev, templated); payload templates; model selection on dev; run plan; pre-registered predictions (Table~\ref{tab:prereg}).}
\begin{table}[h]\centering\small
\begin{tabular}{@{}lll@{}}\toprule
ID & Prediction & Outcome\\\midrule
H1a & G1: M1 $\approx$ M2 (JSON alone does not stop redirection) & \tbd\\
H1b & G2: M2 $<$ M1 & \tbd\\
H2 & G1: M3 $<$ M1, M4 $<$ M2 & \tbd\\
H3 & interaction (M2$-$M4) $>$ (M1$-$M3) & \tbd\\
H4 & gateway: G1 = 0, G2 unchanged & \tbd\\\bottomrule
\end{tabular}
\caption{Pre-registered predictions.}\label{tab:prereg}
\end{table}

% ---------------------------------------------------------------- 0.75 p
\section{Metrics and Analysis}
\todo{ASR-G1/G2, attempt rate, utility (clean, attacked), safe utility, handoff labels, inter-annotator $\kappa$, cost per successful clean task; paired scenario bootstrap.}

% ---------------------------------------------------------------- 1.5 p
\section{Results}
\subsection{Clean utility} \todo{Pipelines must work before safety is compared.}
\subsection{Main effects and interaction} \todo{Fig.\ 2 interaction plot; Table~\ref{tab:main}.}
\begin{table}[h]\centering\small
\begin{tabular}{@{}lccccc@{}}\toprule
Config & Util.\ (clean) & ASR-G1 & ASR-G2 & Util.\ (attacked) & Tokens/ep\\\midrule
A & \tbd & \tbd & \tbd & \tbd & \tbd\\
M1 & \tbd & \tbd & \tbd & \tbd & \tbd\\
M2 & \tbd & \tbd & \tbd & \tbd & \tbd\\
M3 & \tbd & \tbd & \tbd & \tbd & \tbd\\
M4 & \tbd & \tbd & \tbd & \tbd & \tbd\\\bottomrule
\end{tabular}
\caption{Main results (95\% CIs in brackets).}\label{tab:main}
\end{table}
\subsection{Single agent vs.\ pipelines} \todo{overall comparison only}
\subsection{By payload template} \todo{generic vs.\ role-targeted}
\subsection{Gateway} \todo{optional}

% ---------------------------------------------------------------- 1 p
\section{Mechanism Analysis}
\todo{Propagation profile (Fig.\ 3); 2--3 annotated traces; handoff replay (Table 6).}

% ---------------------------------------------------------------- 0.75 p
\section{Discussion and Limitations}
\todo{Implications; one domain, 1--2 models, non-adaptive attacker, synthetic data, model selected for non-zero ASR, prompt sensitivity, judge bias.}

\section{Conclusion}
\todo{Question, one sentence per finding, one design recommendation.}

\bibliographystyle{plainnat}
\bibliography{refs}

\appendix
\section{Handoff schemas}\label{app:schemas}
\todo{Paste from harness/handoff.py.}
\section{Prompts and payload templates}
\todo{Paste from harness/prompts.py and harness/attacks.py.}

\end{document}
